\subsection{结点向量与逆根}

引理 2. 设 S 是 T 的闭子群，Z(S) 是它在 G 中的正规化子。

\begin{enumerate}
\def\labelenumi{(\roman{enumi})}
\item
  \(\mathrm{R}(\mathrm{Z}(\mathrm{S})_0,\mathrm{T})\) 是诸 \(\alpha \in \mathrm{R}(\mathrm{G},\mathrm{T})\) 中满足 \(\alpha (\mathrm{S}) = \{1\}\) 者所成的集合。
\item
  \(\mathrm{Z(S)_0}\) 的中心是诸 \(\operatorname {Ker}\alpha\)（\(\alpha \in \mathrm{R}(\mathrm{Z(S)}_0,\mathrm{T})\)）的交。
\item
  若 S 连通，则 Z(S) 连通。
\end{enumerate}

李代数 \(\mathrm{L}(\mathrm{Z}(\mathrm{S}))_{(\mathbf{C})}\) 由 S 在 \(g_{C}\) 上的不变量组成（第三章，§9，第 3 节，命题 8），因而是 \(t_{C}\) 与诸 \(g^{\alpha}\)（满足 \(\alpha(\mathrm{S})=\{1\}\) 者）的直和，由此得 (i)。断言 (ii) 由命题 8（第 4 节）得出，而断言 (iii) 已在 §2 第 2 节定理 2 的推论 5 中证明。

定理 1. 设 \(\alpha \in \mathrm{R}(\mathrm{G},\mathrm{T})\)。令 \(Z_{\alpha}\) 为 \(\alpha\) 的核的中心化子；它是 \(\mathrm{G}\) 的连通闭子群；其中心是 \(\operatorname{Ker} \alpha\)；其换位子群 \(\mathrm{D}(Z_{\alpha}) = S_{\alpha}\) 是 \(\mathrm{G}\) 的秩 1 连通闭半单子群。我们有 \(\mathrm{R}(Z_{\alpha},\mathrm{T}) = \{\alpha, -\alpha\}\)，且 \(\dim Z_{\alpha} = \dim T + 2\)。

设 \(Z_{\alpha}^{\prime}\) 是 \((\mathrm{Ker}\alpha)_0\) 的中心化子。由引理 2，它是 G 的连通闭子群，且 \(\mathrm{R}(Z_{\alpha}^{\prime},\mathrm{T})\) 是 \(\beta \in \mathrm{R}(\mathrm{G},\mathrm{T})\) 中满足 \(\beta ((\mathrm{Ker}\alpha)_0) = \{1\}\) 者所成的集合。显然，\(\{\alpha , - \alpha \} \subset \mathrm{R}(Z_{\alpha}^{\prime},\mathrm{T})\)。反之，设 \(\beta \in \mathrm{R}(Z_{\alpha}^{\prime},\mathrm{T})\)；由于 \((\mathrm{Ker}\alpha)_0\) 在 \(\mathrm{Ker}\alpha\) 中指数有限，故存在整数 \(r\neq 0\)，使得 \(t^{r\beta} = 1\)（对 \(t\in \mathrm{Ker}\alpha\)）。由序列

\[
0 \longrightarrow \mathbf {Z} \longrightarrow \mathrm{X(T)} \longrightarrow \mathrm{X(Ker} \alpha) \longrightarrow 0
\]

（它通过对偶对应于正合序列

\[
0 \longrightarrow \operatorname{Ker} \alpha \longrightarrow \mathrm{T} \stackrel {{\alpha}} {{\longrightarrow}} \mathrm{U} \longrightarrow 0,
\]

）的正合性可知，\(r\beta\) 是 \(\alpha\) 的倍数；由第八章 §2 第 2 节定理 2 (i)，这蕴含 \(\beta \in \{\alpha, -\alpha\}\)。于是 \(\mathrm{R}(Z_{\alpha}', T) = \{\alpha, -\alpha\}\)。由此并根据引理 2，\(Z_{\alpha}'\) 的中心是 \(\operatorname{Ker} \alpha\)，故 \(Z_{\alpha}' = Z_{\alpha}\)。最后，由 §1 第 4 节命题 4 的推论 1，\(\mathrm{D}(\mathrm{Z}_{\alpha})\) 是 \(\mathrm{G}\) 的连通闭半单子群；它秩为 1，因为 \(\mathcal{DL}(\mathrm{Z}_{\alpha})_{(\mathbf{C})} = \mathfrak{g}^{\alpha} + \mathfrak{g}^{-\alpha} + [\mathfrak{g}^{\alpha},\mathfrak{g}^{-\alpha}]\)。

推论. 存在具有以下性质的李群态射 \(\nu : \mathbf{SU}(2, \mathbf{C}) \to \mathrm{G}\)：

\begin{enumerate}
\def\labelenumi{\alph{enumi})}
\item
  \(\nu\) 的像与 \(\alpha\) 的核交换。
\item
  对所有 \(a \in U\)，我们有 \(\nu\left(\begin{array}{cc}a & 0 \\ 0 & \bar{a}\end{array}\right) \in T\)，且 \(\alpha \circ \nu\left(\begin{array}{cc}a & 0 \\ 0 & \bar{a}\end{array}\right) = a^{2}\)。
\end{enumerate}

若 \(\nu_{1}\) 与 \(\nu_{2}\) 是从 \(\mathbf{SU}(2,\mathbf{C})\) 到 \(\mathrm{G}\) 且具有前述性质的两个态射，则存在 \(a\in \mathbf{U}\) 使得 \(\nu_{2} = \nu_{1}\circ \operatorname {Int}\left( \begin{array}{cc}a & 0\\ 0 & \bar{a} \end{array} \right)\)。

由定理 1 以及 §3 第 6 节命题 6，存在李群态射 \(\nu : \mathbf{SU}(2, \mathbf{C}) \to \mathrm{S}_{\alpha}\)，它是满射且核为离散。于是 \(\nu^{-1}(\mathrm{T} \cap \mathrm{S}_{\alpha})\) 是 \(\mathbf{SU}(2, \mathbf{C})\) 的一个极大环面（§2，第 3 节，命题 1）。由于 \(\mathbf{SU}(2, \mathbf{C})\) 的诸极大环面彼此共轭（§2，第 2 节，定理 2），必要时把 \(\nu\) 换为 \(\nu \circ \operatorname{Int}s\)（其中 \(s \in \mathbf{SU}(2, \mathbf{C})\)），可以假设 \(\nu^{-1}(\mathrm{T} \cap \mathrm{S}_{\alpha})\) 是 \(\mathbf{SU}(2, \mathbf{C})\) 中的对角矩阵群。于是 \(\nu\begin{pmatrix} a & 0 \\ 0 & \bar{a} \end{pmatrix} \in \mathrm{T}\) 对所有 \(a \in \mathbf{U}\) 成立，且映射

\[
\left( \begin{array}{c c} a & 0 \\ 0 & \bar {a} \end{array} \right) \mapsto \alpha \circ \nu \left( \begin{array}{c c} a & 0 \\ 0 & \bar {a} \end{array} \right)
\]

是 \(\mathbf{SU}(2,\mathbf{C})\) 的一个根，从而等于两个映射 \(\begin{pmatrix}a&0\\ 0&\bar{a}\end{pmatrix}\mapsto a^{2}\) 或 \(\begin{pmatrix}a&0\\ 0&\bar{a}\end{pmatrix}\mapsto a^{-2}\) 之一（§3，第 6 节，公式 (19)）。在第一种情形，同态 \(\nu\) 具有所要求的性质；在第二种情形，同态 \(\nu\circ Int\theta\) 具有这些性质（同前处，公式 (18)）。

若 \(\nu_{1}\) 与 \(\nu_{2}\) 是从 \(\mathbf{SU}(2,\mathbf{C})\) 到 G 且满足所述条件的两个态射，则它们都把 \(\mathbf{SU}(2,\mathbf{C})\) 映到 \(S_{\alpha}\) 内（条件 a)），因而都是 \(S_{\alpha}\) 的泛覆盖。于是，存在 \(\varphi\)（\(\mathbf{SU}(2,\mathbf{C})\) 的一个自同构）使得 \(\nu_{2} = \nu_{1} \circ \varphi\)，再应用第 4 节命题 9 即得结论。

由前述推论可知，同态 \(\nu_{\mathrm{T}}\)（从 \(\mathbf{U}\) 到 \(\mathrm{T}\)，由 \(\nu_{\mathrm{T}}(a) = \nu \left( \begin{array}{cc}a & 0\\ 0 & \bar{a} \end{array} \right)\)（\(a\in \mathbf{U}\)）定义）与 \(\nu\) 的选取无关。用 \(K_{\alpha}\in \Gamma (\mathrm{T})\) 表示在 \(\Gamma (\nu_{\mathrm{T}})\) 之下元素 \(2\pi i\)（视为 \(\Gamma (\mathbf{U}) = 2\pi i\mathbf{Z}\) 的元素）所对应的像；它称为相伴于根 \(\alpha\) 的结点向量。我们有 \(\langle \alpha ,K_{\alpha}\rangle = 2\)，换言之（第 2 节，公式 (2)）\(\delta (\alpha)(K_{\alpha}) = 4\pi i\)；由于 \(K_{\alpha}\) 属于 \(\mathfrak{t}\) 与 \(\mathrm{L}(\mathrm{S}_{\alpha})_{(\mathbf{C})}\) 的交，我们有

\[
K _ {\alpha} = 2 \pi i H _ {\delta (\alpha)},\tag{13}
\]

其中 \(H_{\delta(\alpha)}\) 是根 \(\delta(\alpha)\)（\((\mathfrak{g}_{\mathbf{C}}, \mathfrak{t}_{\mathbf{C}})\) 的根）所相伴的逆根（第八章，§2，第 2 节）。换言之，当 \(\Gamma(\mathrm{T}) \otimes \mathbf{R}\) 等同于 \(\mathrm{X(T)}\otimes \mathbf{R}\) 的对偶（借助配对 \(\langle ,\rangle ,K_{\alpha}\)）时，该元素便等同于逆根 \(\alpha^{\vee}\in (\mathrm{X(T)}\otimes \mathbf{R})^{*}\)。

注. 对所有 \(x \in R\)，我们有

\[
\nu \left( \begin{array}{c c} \exp (2 \pi i x) & 0 \\ 0 & \exp (- 2 \pi i x) \end{array} \right) = \nu_ {\mathrm{T}} (e ^ {2 \pi i x}) = \exp (x K _ {\alpha}).\tag{14}
\]

特别地：

\[
\nu \left( \begin{array}{c c} - 1 & 0 \\ 0 & - 1 \end{array} \right) = \nu_ {\mathrm{T}} (- 1) = \exp \Bigl (\frac {1}{2} K _ {\alpha} \Bigr).\tag{15}
\]

由此可知，\(\nu\) 是单射当且仅当 \(K_{\alpha} \notin 2\Gamma(T)\)，换言之，当且仅当存在 \(\lambda \in \mathrm{X}(T)\) 使得 \(\langle \lambda, K_{\alpha} \rangle \notin 2Z\)。当 \(g_{C}\) 为单时，\(\nu\) 都是单射，除非 \(g_{C}\) 是 \(B_{n}\) 型、\(C(G) = \{1\}\) 且 \(\alpha\) 是短根（参看第六章，图版）。

在本段余下部分，我们用 \(\mathrm{R}^{\vee}(\mathrm{G},\mathrm{T})\) 表示诸结点向量 \(K_{\alpha}\)（\(\alpha\in\mathrm{R}(\mathrm{G},\mathrm{T})\)）所成的集合。它是 \(\varGamma(\mathrm{T})\) 的一个子集，\(\varGamma(\mathrm{T})\) 到 \(t_{C}\) 中的典则单射把它等同于对比率为 \(2\pi i\) 的位似作用于逆根系 \(\mathrm{R}^{\vee}(\mathfrak{g}_{\mathbf{C}},\mathfrak{t}_{\mathbf{C}})=\{H_{\delta(\alpha)}\}\)（即 \(\delta(\mathrm{R})\) 的逆根系）所得的集合。由此可知，\(\mathrm{R}^{\vee}(\mathrm{G},\mathrm{T})\) 生成 R-向量空间 \(\mathrm{L}(\mathrm{T}\cap\mathrm{D}(\mathrm{G}))\)，从而它在 \(\mathrm{X}(\mathrm{T})\) 中的正交补是 \(\mathrm{X}(\mathrm{T}/(\mathrm{T}\cap\mathrm{D}(\mathrm{G})))\)。

用 \(\mathrm{Aut(T)}\) 表示李群 T 的自同构群；Weyl 群 \(W = W_{G}(T)\)（§2，第 5 节）可等同于 \(\mathrm{Aut(T)}\) 的一个子群。另一方面，回忆（第八章，§2，第 2 节，注 4）：Weyl 群 \(W(\mathfrak{g}_{\mathbf{C}}, t_{\mathbf{C}})\)（分裂约化代数（\(\mathfrak{g}_{\mathbf{C}}, t_{\mathbf{C}}\)）的 Weyl 群）作用在 \(t_{\mathbf{C}}\) 上，从而典则地等同于 \(\mathbf{GL}(t_{\mathbf{C}})\) 的一个子群。

命题 10. 映射 \(u \mapsto \mathrm{L}(u)_{(\mathbf{C})}\)（从 \(\operatorname{Aut}(\mathrm{T})\) 到 \(\mathbf{GL}(\mathfrak{t}_{\mathbf{C}})\)）诱导出从 \(\mathrm{W}\) 到分裂约化李代数 \((\mathfrak{g}_{\mathbf{C}}, \mathfrak{t}_{\mathbf{C}})\) 的 Weyl 群的一个同构。对所有 \(\alpha \in \mathrm{R}\)，\(\mathrm{W}_{\mathrm{Z}_{\alpha}}(\mathrm{T})\) 的阶为 2，且 \(\mathrm{W}_{\mathrm{Z}_{\alpha}}(\mathrm{T})\) 的非恒等元在上述同构下的像是反射 \(s_{H_{\delta(\alpha)}}\)。

所考虑的映射是单射。剩下要证明它的像等于 \(\mathrm{W}(\mathfrak{g}_{\mathbf{C}}, t_{\mathbf{C}})\)。

设 \(g \in \mathrm{N}_{\mathrm{G}}(\mathrm{T})\)。采用第八章 §5 第 2 节的记号，我们有 \(\operatorname{Ad} g \in \operatorname{Aut}(\mathfrak{g}_{\mathbf{C}}, \mathfrak{t}_{\mathbf{C}}) \cap \operatorname{Int}(\mathfrak{g}_{\mathbf{C}})\)，故 \(\operatorname{Ad} g \in \operatorname{Aut}_{0}(\mathfrak{g}_{\mathbf{C}}, \mathfrak{t}_{\mathbf{C}})\)（同前处，第 5 节，命题 11）。由同前处第 2 节命题 4，\(\mathfrak{t}_{\mathbf{C}}\) 上由 \(\operatorname{Ad} g\) 诱导的自同构属于 \(\mathrm{W}(\mathfrak{g}_{\mathbf{C}}, \mathfrak{t}_{\mathbf{C}})\)。于是，\(\mathrm{W}\) 在 \(\mathbf{GL}(\mathfrak{t}_{\mathbf{C}})\) 中的像含于 \(\mathrm{W}(\mathfrak{g}_{\mathbf{C}}, \mathfrak{t}_{\mathbf{C}})\)。

设 \(\alpha \in \mathrm{R}(\mathrm{G},\mathrm{T})\)，并设 \(\nu : \mathbf{SU}(2,\mathbf{C}) \to \mathrm{G}\) 是具有定理 1 的推论中所述性质的李群态射。\(\nu\) 在元素 \(\theta\)（\(\mathbf{SU}(2,\mathbf{C})\) 的元素）上的像具有以下性质（§3，第 6 节，公式 (17)）：

\[
a) (\operatorname{Int} \nu (\theta)) (t) = t \quad \text{ 若 } t \in \operatorname{Ker} \alpha ,
\]

\[
b) (\operatorname{Int} \nu (\theta)) (t) = t ^ {- 1} \quad \text{ 若 } t \in \mathrm{T} \cap \mathrm{S} _ {\alpha}.
\]

由此可知，\(\operatorname{Ad}\nu(\theta)\) 在 \(\operatorname{Ker}\delta(\alpha)\subset\mathfrak{t}_{\mathbf{C}}\) 上诱导出恒等映射，并诱导出映射 \(x\mapsto-x\)（在 \([\mathfrak{g}^{\alpha},\mathfrak{g}^{-\alpha}]\) 上），因而与反射 \(s_{H_{\delta(\alpha)}}\) 一致。于是，W 的像包含所有的 \(s_{H_{\delta(\alpha)}}\)，从而等于 \(\mathrm{W}(\mathfrak{g}_{\mathbf{C}}, t_{\mathbf{C}})\)。特别地，\(\mathrm{W}_{\mathrm{Z}_{\alpha}}(\mathrm{T})\) 的阶为 2，因而由恒等映射与 \(\operatorname{Int}\nu(\theta)\) 组成。命题证毕。

推论. 假设 G 是半单的。则 G 的每个元素都是 G 中两个元素的换位子。

设 \(c\) 是 Weyl 群 \(\mathrm{W}(\mathfrak{g}_{\mathbf{C}},\mathfrak{t}_{\mathbf{C}})\) 的一个 Coxeter 变换（第五章，§6，第 1 节），并设 \(n\) 是 \(\mathrm{N}_{\mathrm{G}}(\mathrm{T})\) 的一个元素，它在 \(\mathrm{W}\) 中的类经由命题中定义的同构等同于 \(c\)。用 \(f_{c}\) 表示态射 \(t\mapsto (n,t)\)（从 \(\mathrm{T}\) 到 \(\mathrm{T}\)）；对 \(x\in \mathfrak{t}_{\mathbf{C}}\)，我们有 \(\mathrm{L}(f_c)_{(\mathbf{C})}(x) = (\mathrm{Ad}n)(x) - x = c(x) - x\)。

由第五章 §6 第 2 节定理 1，自同态 \(c\)（\(\mathbf{t}_{\mathbf{C}}\) 上的）没有等于 1 的特征值。因此，\(\mathrm{L}(f_c)\) 是满射，从而 \(f_c\) 也是满射。由此可知，T 的每个元素都是 G 中两个元素的换位子，结合 §2 第 2 节定理 2 即得推论。

